% GENERATED FILE. Edit canonical lessons, then run npm run generate:books.
% canonical-source-hash: fnv1a64:1a37d375fe14e17d
% canonical-lessons: AR-C194-ibtalaa, AR-C194-madagha, AR-C194-lamasa, AR-C194-shamma, AR-C194-amsaka

\chapter{To swallow, To chew, To touch, To smell, To hold}
\label{ch:ar-to-swallow-to-chew-to-touch-to-smell-to-hold}
\hlchaptermodality{194}

\begin{chapteropening}
\textbf{By the end of this chapter:} \emph{I can say to swallow, to chew, to touch, to smell and to hold.}

Complete the last lesson of chapter 194: I can say to swallow, to chew, to touch, to smell and to hold.
\end{chapteropening}

\section[\texorpdfstring{\ar{ابتلع} (ibtalaʿa)}{ابتلع (ibtalaʿa)}]{\ar{ابتلع} (ibtalaʿa) — to swallow}

\label{lesson:AR-C194-ibtalaa}



\begin{quote}
Before the new one: say the Arabic for to bargain, then the Arabic for to go shopping.
\end{quote}

\subsection*{You'll want to know: \ar{ابتلع}}
\textbf{\ar{ابتلع}} — \emph{ibtalaʿa} — \textquotedblleft{}to swallow\textquotedblright{}.

\begin{cousinweb}
\textbf{\ar{ابتلع}} (\emph{ibtalaʿa}) is built on the native Arabic root \emph{b-l-ʿ}, \textquotedblleft{}to swallow\textquotedblright{} (Lane). \emph{Ibtalaʿa} is Form VIII, \textquotedblleft{}to swallow\textquotedblright{}.
\end{cousinweb}

\begin{grammarlens}[title={What you've built}]
One more word for this chapter.
\end{grammarlens}

\subsection*{Your turn}
\begin{itemize}
\raggedright
  \item \emph{Say it:} \emph{ibtalaʿa}
  \item \emph{Say it:} \emph{ibtalaʿa}, once more
  \item \emph{Say it:} say \emph{ibtalaʿa}
  \item \emph{Recall:} say the Arabic for to bargain, then the Arabic for to go shopping, then say \emph{ibtalaʿa} again
\end{itemize}

\begin{tcolorbox}[breakable,colback=teal!4,colframe=teal!35!black,title={Before you move on}]
What does \ar{ابتلع} mean? (\textquotedblleft{}To swallow\textquotedblright{}.) Say it once more.
\end{tcolorbox}
\section[\texorpdfstring{\ar{مضغ} (maḍagha)}{مضغ (maḍagha)}]{\ar{مضغ} (maḍagha) — to chew}

\label{lesson:AR-C194-madagha}



\begin{quote}
Before the new one: say the Arabic for to go shopping, then the Arabic for to swallow.
\end{quote}

\subsection*{You'll want to know: \ar{مضغ}}
\textbf{\ar{مضغ}} — \emph{maḍagha} — \textquotedblleft{}to chew\textquotedblright{}.

\begin{cousinweb}
\textbf{\ar{مضغ}} (\emph{maḍagha}) is built on the native Arabic root \emph{m-ḍ-gh}, \textquotedblleft{}to chew\textquotedblright{} (Lane). \emph{Maḍagha} is the plain verb, \textquotedblleft{}he chewed\textquotedblright{}.
\end{cousinweb}

\begin{grammarlens}[title={What you've built}]
One more word for this chapter.
\end{grammarlens}

\subsection*{Your turn}
\begin{itemize}
\raggedright
  \item \emph{Say it:} \emph{maḍagha}
  \item \emph{Say it:} \emph{maḍagha}, once more
  \item \emph{Say it:} say \emph{maḍagha}
  \item \emph{Recall:} say the Arabic for to go shopping, then the Arabic for to swallow, then say \emph{maḍagha} again
\end{itemize}

\begin{tcolorbox}[breakable,colback=teal!4,colframe=teal!35!black,title={Before you move on}]
What does \ar{مضغ} mean? (\textquotedblleft{}To chew\textquotedblright{}.) Say it once more.
\end{tcolorbox}
\section[\texorpdfstring{\ar{لمس} (lamasa)}{لمس (lamasa)}]{\ar{لمس} (lamasa) — to touch}

\label{lesson:AR-C194-lamasa}



\begin{quote}
Before the new one: say the Arabic for to swallow, then the Arabic for to chew.
\end{quote}

\subsection*{You'll want to know: \ar{لمس}}
\textbf{\ar{لمس}} — \emph{lamasa} — \textquotedblleft{}to touch\textquotedblright{}.

\begin{cousinweb}
\textbf{\ar{لمس}} (\emph{lamasa}) is built on the native Arabic root \emph{l-m-s}, \textquotedblleft{}to touch\textquotedblright{} (Lane). \emph{Lamasa} is the plain verb, \textquotedblleft{}he touched\textquotedblright{}.
\end{cousinweb}

\begin{grammarlens}[title={What you've built}]
One more word for this chapter.
\end{grammarlens}

\subsection*{Your turn}
\begin{itemize}
\raggedright
  \item \emph{Say it:} \emph{lamasa}
  \item \emph{Say it:} \emph{lamasa}, once more
  \item \emph{Say it:} say \emph{lamasa}
  \item \emph{Recall:} say the Arabic for to swallow, then the Arabic for to chew, then say \emph{lamasa} again
\end{itemize}

\begin{tcolorbox}[breakable,colback=teal!4,colframe=teal!35!black,title={Before you move on}]
What does \ar{لمس} mean? (\textquotedblleft{}To touch\textquotedblright{}.) Say it once more.
\end{tcolorbox}
\section[\texorpdfstring{\ar{شم} (shamma)}{شم (shamma)}]{\ar{شم} (shamma) — to smell}

\label{lesson:AR-C194-shamma}



\begin{quote}
Before the new one: say the Arabic for to chew, then the Arabic for to touch.
\end{quote}

\subsection*{You'll want to know: \ar{شم}}
\textbf{\ar{شم}} — \emph{shamma} — \textquotedblleft{}to smell\textquotedblright{}.

\begin{cousinweb}
\textbf{\ar{شم}} (\emph{shamma}) is built on the native Arabic root \emph{sh-m-m}, \textquotedblleft{}to smell\textquotedblright{} (Lane). \emph{Shamma} is the plain verb, \textquotedblleft{}he smelled\textquotedblright{}.
\end{cousinweb}

\begin{grammarlens}[title={What you've built}]
One more word for this chapter.
\end{grammarlens}

\subsection*{Your turn}
\begin{itemize}
\raggedright
  \item \emph{Say it:} \emph{shamma}
  \item \emph{Say it:} \emph{shamma}, once more
  \item \emph{Say it:} say \emph{shamma}
  \item \emph{Recall:} say the Arabic for to chew, then the Arabic for to touch, then say \emph{shamma} again
\end{itemize}

\begin{tcolorbox}[breakable,colback=teal!4,colframe=teal!35!black,title={Before you move on}]
What does \ar{شم} mean? (\textquotedblleft{}To smell\textquotedblright{}.) Say it once more.
\end{tcolorbox}
\section[\texorpdfstring{\ar{أمسك} (amsaka)}{أمسك (amsaka)}]{\ar{أمسك} (amsaka) — to hold, to grab}

\label{lesson:AR-C194-amsaka}



\begin{quote}
Before the new one: say the Arabic for to touch, then the Arabic for to smell.
\end{quote}

\subsection*{You'll want to know: \ar{أمسك}}
\textbf{\ar{أمسك}} — \emph{amsaka} — \textquotedblleft{}to hold, to grab\textquotedblright{}.

\begin{cousinweb}
\textbf{\ar{أمسك}} (\emph{amsaka}) is built on the native Arabic root \emph{m-s-k}, \textquotedblleft{}to hold, to grasp\textquotedblright{} (Lane). \emph{Amsaka} is Form IV, \textquotedblleft{}to take hold\textquotedblright{}.
\end{cousinweb}

\begin{grammarlens}[title={What you've built}]
One more word for this chapter.
\end{grammarlens}

\subsection*{Your turn}
\begin{itemize}
\raggedright
  \item \emph{Say it:} \emph{amsaka}
  \item \emph{Say it:} \emph{amsaka}, once more
  \item \emph{Say it:} say \emph{amsaka}
  \item \emph{Recall:} say the Arabic for to touch, then the Arabic for to smell, then say \emph{amsaka} again
\end{itemize}

\begin{tcolorbox}[breakable,colback=teal!4,colframe=teal!35!black,title={Before you move on}]
What does \ar{أمسك} mean? (\textquotedblleft{}To hold, to grab\textquotedblright{}.) Say it once more.
\end{tcolorbox}
